\documentclass[conference]{IEEEtran}
\IEEEoverridecommandlockouts
\usepackage{cite}
\usepackage{amsmath,amssymb,amsfonts}
\usepackage{algorithmic}
\usepackage{graphicx}
\usepackage{textcomp}
\usepackage{xcolor}
\usepackage{multirow}
\usepackage[brazil]{babel}
\def\BibTeX{{\rm B\kern-.05em{\sc i\kern-.025em b}\kern-.08em
    T\kern-.1667em\lower.7ex\hbox{E}\kern-.125emX}}

\begin{document}

\title{Análise Comparativa de Arquiteturas de Deep Learning para
Classificação de Retinopatia Diabética}

\author{
\IEEEauthorblockN{
Gabriel de Sousa Nobre,
Victor Lins Gurgel do Amaral,
Lorenzo Barros Calheiros Pinheiro,
Davi Oliveira de Alencar}
\IEEEauthorblockA{
\textit{Graduação em Ciência da Computação} \\
\textit{Universidade de Fortaleza (UNIFOR)} \\
Fortaleza, Ceará, Brasil \\
gabriel@email.com,
vlinsgurgel@email.com,
lorenzo@email.com,
davi@email.com}}

\maketitle

\begin{abstract}
A retinopatia diabética é uma importante causa de perda de visão,
tornando relevante o desenvolvimento de métodos automatizados para
apoiar sua triagem e diagnóstico. Este trabalho avaliou diferentes
arquiteturas de aprendizado profundo para a classificação de
imagens de retina em cinco níveis de severidade da doença,
utilizando o dataset \textit{Diabetic Retinopathy (resized)}. Foram
comparados modelos baseados em redes neurais convolucionais e Vision
Transformers, empregando transferência de aprendizado, fine-tuning,
aumento de dados e ponderação de classes para lidar com o
desbalanceamento da base. Os resultados indicaram que arquiteturas
modernas, como ConvNeXt-Small e EfficientNetV2-S, apresentaram os
melhores desempenhos gerais. O ConvNeXt-Small obteve o maior coeficiente
Kappa, enquanto o EfficientNetV2-S apresentou melhor equilíbrio entre
Kappa e F1 macro. Também foi avaliada uma estratégia de ensemble
baseada em \textit{Weighted Soft Voting}; entretanto, essa abordagem
não superou os modelos individuais na configuração testada. Os
resultados sugerem que a resolução das imagens, o ajuste fino das
camadas e o tratamento do desbalanceamento são fatores determinantes
para o desempenho na classificação multiclasse de retinopatia
diabética.
\end{abstract}

\begin{IEEEkeywords}
Retinopatia diabética, classificação de imagens médicas,
Deep Learning, CNNs, Vision Transformers, Transfer Learning.
\end{IEEEkeywords}

% =============================================================
\section{Introdução}
% =============================================================

A retinopatia diabética (RD) é uma complicação microvascular
do diabetes mellitus que afeta os vasos sanguíneos da retina e
constitui uma das principais causas de cegueira evitável entre adultos
em idade produtiva. O crescimento contínuo da população
diabética mundial tem aumentado significativamente a demanda por
programas de rastreamento e detecção precoce da doença.
Segundo Grzybowski \cite{grzybowski2022}, a expansão global do
diabetes e o aumento da incidência de retinopatia diabética tornam
cada vez mais necessária a adoção de soluções escaláveis
para triagem oftalmólogica, especialmente em regiões com recursos
limitados e escassez de especialistas.

Tradicionalmente, a identificação da retinopatia diabética é
realizada por oftalmologistas por meio da análise de imagens de
retinografia. Entretanto, a grande quantidade de exames gerados em
programas de rastreamento populacional torna esse processo demorado,
oneroso e dependente da disponibilidade de profissionais qualificados.
Nesse contexto, técnicas de inteligência artificial, especialmente
aquelas baseadas em aprendizado profundo, têm demonstrado grande
potencial para automatizar a análise de imagens médicas e auxiliar
a tomada de decisão clínica.

Entre os trabalhos pioneiros da área, Gulshan et al. \cite{gulshan2016}
desenvolveram um sistema baseado em redes neurais convolucionais profundas
para detecção automática de retinopatia diabética em fotografias
de retina. O modelo alcançou elevados índices de sensibilidade e
especificidade em bases de validação independentes, demonstrando
que algoritmos de aprendizado profundo podem atingir desempenho
comparável ao de especialistas em tarefas de triagem da doença.

Nos últimos anos, avanços arquiteturais têm impulsionado ainda
mais o desempenho desses sistemas. Além das tradicionais Redes Neurais
Convolucionais (Convolutional Neural Networks -- CNNs), arquiteturas
baseadas em mecanismos de atenção, conhecidas como Vision
Transformers (ViTs), passaram a apresentar resultados competitivos em
diversas tarefas de visão computacional. Em um estudo comparativo
envolvendo diferentes doenças da retina, Sayg{\i}l{\i} e At{\i}c{\i}
\cite{saygili2025} observaram que o modelo Swin Transformer superou
arquiteturas CNN convencionais e outros modelos Transformer avaliados,
alcançando a maior acurácia entre os métodos analisados.

Paralelamente, abordagens que combinam múltiplos modelos também vêm
sendo investigadas com o objetivo de aumentar a robustez das
predições. Rajeshwar et al. \cite{rajeshwar2025} propuseram uma
arquitetura híbrida baseada em extração de características por
aprendizado profundo e ensemble por stacking, obtendo resultados
expressivos na detecção de retinopatia diabética. Esses estudos
evidenciam o interesse crescente da comunidade científica em explorar
diferentes arquiteturas e estratégias de combinação de modelos
para melhorar o desempenho dos sistemas automatizados de apoio \`a
triagem oftalmólogica.

Apesar dos avanços observados na literatura, a classificação
automática da retinopatia diabética ainda apresenta desafios
relevantes. Um deles é o desbalanceamento entre as classes, uma vez
que bases de dados públicas frequentemente possuem predomin\^ancia de
imagens sem sinais da doença em comparação com os níveis
mais severos. Além disso, trata-se de uma tarefa multiclasse e
ordinal, pois os níveis de severidade seguem uma progressão clínica
de 0 a 4, como pode ser observado na Fig.~\ref{DR_levels}. Dessa forma,
métricas globais como acurácia podem não refletir adequadamente o
desempenho do modelo em classes minoritárias, tornando necessária a
análise de métricas como F1 macro e coeficiente Kappa ponderado.

\begin{figure*}[htbp]
\centering
\includegraphics[width=0.9\textwidth]{DR_levels.png}
\caption{Imagens do dataset por nível de severidade.}
\label{DR_levels}
\end{figure*}

Diante desse contexto, este trabalho tem como objetivo avaliar e comparar
diferentes arquiteturas de aprendizado profundo para a classificação
de retinopatia diabética em imagens do fundo da retina. Foram
investigados modelos baseados em CNNs e Vision Transformers, incluindo
ResNet50, DenseNet121, DenseNet169, EfficientNet-B0, EfficientNet-B2,
EfficientNetV2-S, ConvNeXt-Small e Swin-Tiny. Também foram avaliadas
diferentes estratégias de fine-tuning, aumento de dados, balanceamento
por pesos de classe e ensemble, buscando identificar as configurações
mais adequadas para a tarefa proposta.

% =============================================================
\section{Metodologia}
% =============================================================

A metodologia adotada consistiu na comparação experimental de
diferentes arquiteturas de aprendizado profundo para a classificação
multiclasse de retinopatia diabética em imagens de fundo de retina.
O fluxo experimental envolveu a preparação da base de dados, o
pré-processamento das imagens, o treinamento de modelos pré-treinados
por transferência de aprendizado, a aplicação de estratégias
de balanceamento e a avaliação dos modelos por métricas
adequadas a um problema ordinal e desbalanceado.

\subsection{Dataset}

Foi utilizado o dataset público \textit{Diabetic Retinopathy
(resized)}\cite{dataset}, derivado da competição de retinopatia
diabética do Kaggle \cite{APTOS2019}. A base contém 35.126 imagens
de fundo de retina, rotuladas em cinco níveis de severidade: classe 0,
sem retinopatia diabética; classe 1, retinopatia leve; classe 2,
retinopatia moderada; classe 3, retinopatia severa; e classe 4,
retinopatia proliferativa. Como esses rótulos representam níveis
progressivos de gravidade clínica, o problema foi tratado como uma
tarefa de classificação multiclasse ordinal.

Os metadados originais da base são disponibilizados em um arquivo CSV
contendo o identificador da imagem e o respectivo nível de severidade.
Para otimizar o treinamento dos modelos, as imagens foram reorganizadas
em diretórios separados por classe, utilizando uma estrutura compatível
com o carregamento por \textit{ImageFolder}. O conjunto foi dividido em
treino e validação na proporção 80/20, com separação
estratificada por classe e semente fixa igual a 42. Não foi utilizado
um conjunto de teste externo independente.

\begin{table}[htbp]
\caption{Distribuição da base de dados por classe}
\label{tab:distribuicao_base}
\centering
\scriptsize
\begin{tabular}{c l r r r}
\hline
\textbf{Classe} & \textbf{Descrição} & \textbf{Total} &
\textbf{Treino} & \textbf{Val.} \\
\hline
0 & No DR         & 25.810 & 20.648 & 5.162 \\
1 & Mild          &  2.443 &  1.955 &   488 \\
2 & Moderate      &  5.292 &  4.234 & 1.058 \\
3 & Severe        &    873 &    699 &   174 \\
4 & Proliferative &    708 &    567 &   141 \\
\hline
\textbf{Total} & -- & \textbf{35.126} & \textbf{28.103} & \textbf{7.023} \\
\hline
\end{tabular}
\end{table}

Observa-se forte desbalanceamento entre as classes, com predomin\^ancia
da classe 0, correspondente a imagens sem sinais de retinopatia diabética.
Esse fator motivou o uso de métricas além da acurácia e a aplicação
de estratégias de ponderação durante o treinamento.

\subsection{Pré-processamento e balanceamento dos dados}

As imagens foram redimensionadas de acordo com a configuração
experimental de cada modelo, sendo avaliadas resoluções de
224$\times$224, 256$\times$256, 384$\times$384, 512$\times$512 e
1024$\times$1024 pixels. Também foi aplicada normalização com
média e desvio padrão do ImageNet, uma vez que os modelos utilizados
foram inicializados com pesos pré-treinados nesse conjunto.

Para aumentar a variabilidade dos dados de treinamento e reduzir o risco
de sobreajuste, foram aplicadas técnicas de \textit{data augmentation},
incluindo inversões horizontal e vertical, rotações de até
15 graus e alterações de cor por meio de \textit{ColorJitter}.
Essas transformações foram aplicadas apenas ao conjunto de
treinamento, preservando o conjunto de validação sem alterações
artificiais além do redimensionamento e da normalização.

Devido ao desbalanceamento observado entre as classes, foi utilizada a
função de perda \textit{CrossEntropyLoss} ponderada por pesos de
classe, calculados a partir da distribuição dos rótulos no
conjunto de treinamento. Estratégias adicionais, como
\textit{WeightedRandomSampler} e \textit{Focal Loss}, também foram
exploradas, mas os principais experimentos utilizaram predominantemente
a ponderação da função de perda.

\subsection{Arquiteturas avaliadas e variações experimentais}

Foram avaliadas arquiteturas baseadas em Redes Neurais Convolucionais e
Vision Transformers. Entre os modelos convolucionais, foram testadas as
arquiteturas EfficientNet-B0, EfficientNet-B2, EfficientNetV2-S,
ResNet50, DenseNet121, DenseNet169 e ConvNeXt-Small. Como representante
baseado em mecanismos de atenção visual, foi avaliado o modelo
Swin-Tiny.

Todos os modelos foram utilizados com pesos pré-treinados em ImageNet,
caracterizando uma abordagem de transferência de aprendizado. A camada
final de classificação foi substituída para produzir cinco
saídas, correspondentes aos níveis de severidade da retinopatia
diabética. As variações experimentais envolveram diferentes
resoluções de entrada, graus de congelamento do \textit{backbone},
otimizadores, taxas de aprendizado, tamanhos de lote e estratégias de
fine-tuning.

Inicialmente, foram avaliadas arquiteturas com maior nível de
congelamento das camadas pré-treinadas, funcionando como
configurações de referência. Em seguida, os modelos com
resultados mais promissores foram submetidos a ajustes mais específicos,
incluindo descongelamento parcial ou total do \textit{backbone}, aumento
da resolução de entrada e alteração do otimizador. Essa
estratégia permitiu analisar o impacto do fine-tuning e da
resolução das imagens no desempenho dos classificadores.

\subsection{Estratégia de treinamento}

Os experimentos foram implementados em Python utilizando as bibliotecas
PyTorch e torchvision, executados em GPUs NVIDIA compatíveis com CUDA.
Em razão das limitações de memória gráfica disponível,
os tamanhos de lote foram ajustados conforme a arquitetura e a
resolução das imagens utilizadas.

Foram utilizados os otimizadores Adam e AdamW, dependendo da
configuração experimental. Em parte dos experimentos, foram
aplicadas taxas de aprendizado distintas para o \textit{backbone} e
para a cabeça classificadora, permitindo ajustes mais suaves nas
camadas pré-treinadas e maior adaptação da camada final. Também
foi utilizado o agendador \textit{ReduceLROnPlateau}, reduzindo a taxa de
aprendizado quando não havia melhora na métrica de validação.

O número máximo de épocas variou entre 25 e 50, conforme o modelo
avaliado. Para evitar sobreajuste, foi empregado \textit{early stopping},
monitorando principalmente o coeficiente Kappa ponderado quadrático no
conjunto de validação. A versão final considerada para cada
modelo correspondeu ao melhor \textit{checkpoint} obtido durante o
treinamento, e não necessariamente \`a última época executada.

\subsection{Estratégia de classificação e ensemble}

A tarefa foi formulada como classificação multiclasse direta. Para
cada imagem de entrada, o modelo produzia uma distribuição de
probabilidades sobre as cinco classes, sendo atribuída como
predição final a classe com maior probabilidade.

Além da avaliação individual dos modelos, foi testada uma
estratégia de ensemble baseada em \textit{Weighted Soft Voting}. Nessa
abordagem, as probabilidades previstas por diferentes modelos foram
combinadas por média ponderada, utilizando pesos proporcionais ao
desempenho individual de cada arquitetura no conjunto de validação.
O ensemble foi avaliado com o objetivo de verificar se a combinação
de modelos de famílias distintas poderia melhorar a robustez das
predições em comparação aos classificadores individuais.

\subsection{Métricas de avaliação}

O desempenho dos modelos foi avaliado por meio de acurácia,
precisão macro, recall macro, F1-score macro e coeficiente Kappa
ponderado quadrático. A acurácia foi utilizada como medida geral da
taxa de acerto; entretanto, devido ao forte desbalanceamento entre as
classes, essa métrica não foi considerada suficiente de forma
isolada.

O F1-score macro foi utilizado por atribuir o mesmo peso a todas as
classes, independentemente da frequência de cada uma na base, o que
permite avaliar melhor o desempenho dos modelos nas classes
minoritárias. O coeficiente Kappa ponderado quadrático foi adotado
como métrica principal de comparação, pois considera a natureza
ordinal do problema, penalizando de forma mais intensa erros entre
classes distantes na escala de severidade. Também foram analisadas
matrizes de confusão para identificar os principais padrões de erro
entre os níveis de retinopatia diabética.

% =============================================================
\section{Resultados e Discussão}
% =============================================================

\subsection{Comparação geral entre os modelos}

A Tabela~\ref{tab:comparacao_modelos} reúne os principais resultados
obtidos. Para não sobrecarregar a comparação com todos os
experimentos realizados, foram incluídos o baseline, o melhor
representante de cada família arquitetural e o ensemble final.

\begin{table}[htbp]
\caption{Comparação dos principais modelos avaliados}
\label{tab:comparacao_modelos}
\centering
\scriptsize
\begin{tabular}{l c c c c}
\hline
\textbf{Modelo} & \textbf{Res.} & \textbf{Accuracy} & \textbf{F1 Macro} & \textbf{Kappa} \\
\hline
EfficientNet-B0 (Baseline) & 224  & 0,74 & 0,27 & 0,26 \\
Swin-Tiny                  & 224  & 0,69 & 0,52 & 0,63 \\
DenseNet169                & 512  & 0,75 & 0,55 & 0,71 \\
EfficientNetV2-S           & 512  & 0,80 & 0,54 & 0,73 \\
ConvNeXt-Small             & 1024 & 0,82 & 0,53 & 0,75 \\
Ensemble                   & --   & 0,73 & 0,39 & 0,47 \\
\hline
\end{tabular}
\end{table}

O caso do EfficientNet-B0 é ilustrativo de um problema recorrente em
bases desbalanceadas: acurácia de 0,74 com kappa de apenas 0,26
indica que o modelo aprendeu, na prática, a prever quase sempre a
classe \textit{No\,DR} sem real capacidade de distinguir os níveis da
doença. Essa discrep\^ancia reforça a import\^ancia de avaliar
os modelos pelo kappa quadrático e pelo F1 macro, e não apenas pela
acurácia global.

\subsection{Melhores modelos individuais}

Os três modelos com melhor desempenho --- ConvNeXt-Small,
EfficientNetV2-S e DenseNet169 --- apresentaram perfis bastante
distintos, e essa diferença é relevante do ponto de vista
clínico.

O ConvNeXt-Small obteve o maior kappa do estudo (0,7453), mas com uma
fragilidade importante: recall de apenas 0,05 para a classe
\textit{Mild} (F1~=~0,09). Na prática, isso significa que o modelo
quase não detecta o estágio inicial da retinopatia, que é
justamente onde a intervenção precoce teria maior impacto. O bom
resultado geral acaba mascarado pela domin\^ancia da classe
\textit{No\,DR}, que corresponde a 73\% das amostras e teve recall de
0,93.

O EfficientNetV2-S em 512\,px mostrou o melhor equilíbrio geral entre
as cinco classes, com kappa de 0,7313 e F1 macro de 0,54. Os recalls de
0,42 para \textit{Severe} e 0,67 para \textit{Proliferative} indicam
que o modelo conseguiu aprender características das classes mais
graves, mesmo com muito menos exemplos disponíveis. A classe
\textit{Mild} permaneceu a mais difícil (recall 0,34; F1 0,27) ---
um resultado esperado, considerando que essa classe tem cerca de 10
vezes menos amostras que \textit{No\,DR} e que suas alterações
vasculares são visualmente sutis, podendo se confundir facilmente com
retinas saudáveis.

O DenseNet169 em 512\,px fechou o grupo dos melhores com kappa 0,7114 e
F1 macro de 0,55. Sua convergência mais lenta --- pico na
época\,18 --- sugere um treinamento mais gradual e estável, e seu
ponto forte ficou no recall de \textit{Proliferative} (0,67), o melhor
entre todos os modelos testados. A Tabela~\ref{tab:metricas_classe}
detalha as métricas por classe dos três modelos.

\begin{table}[htbp]
\caption{Métricas por classe --- melhores modelos individuais}
\label{tab:metricas_classe}
\centering
\scriptsize
\setlength{\tabcolsep}{3.5pt}
\begin{tabular}{l l c c c}
\hline
\textbf{Modelo} & \textbf{Classe} & \textbf{Prec.} & \textbf{Rec.} & \textbf{F1} \\
\hline
\multirow{5}{*}{\shortstack[l]{ConvNeXt-S \\ (1024\,px)}}
  & No DR    & 0,90 & 0,93 & 0,92 \\
  & Mild     & 0,41 & 0,05 & 0,09 \\
  & Moderate & 0,55 & 0,48 & 0,51 \\
  & Severe   & 0,52 & 0,28 & 0,36 \\
  & Prolif.  & 0,71 & 0,55 & 0,62 \\
\hline
\multirow{5}{*}{\shortstack[l]{EffNetV2-S \\ (512\,px)}}
  & No DR    & 0,90 & 0,85 & 0,88 \\
  & Mild     & 0,22 & 0,34 & 0,27 \\
  & Moderate & 0,59 & 0,60 & 0,60 \\
  & Severe   & 0,44 & 0,42 & 0,43 \\
  & Prolif.  & 0,57 & 0,67 & 0,61 \\
\hline
\multirow{5}{*}{\shortstack[l]{DenseNet169 \\ (512\,px)}}
  & No DR    & 0,89 & 0,85 & 0,87 \\
  & Mild     & 0,17 & 0,30 & 0,22 \\
  & Moderate & 0,62 & 0,51 & 0,56 \\
  & Severe   & 0,39 & 0,53 & 0,45 \\
  & Prolif.  & 0,67 & 0,67 & 0,67 \\
\hline
\end{tabular}
\end{table}

\subsection{Curvas de aprendizagem}

% Descomente este bloco quando a figura estiver pronta:
% \begin{figure}[htbp]
% \centering
% \includegraphics[width=\columnwidth]{curvas_aprendizado.png}
% \caption{Curvas de \textit{loss} e kappa quadrático de
% validação por época para os três melhores modelos.}
% \label{fig:curvas}
% \end{figure}

As curvas de convergência mostram comportamentos bem diferentes entre
os modelos. Na Fig.\ref{fig:acc_v2} o EfficientNetV2-S convergiu de forma estável ao longo de
14 épocas, sem sinais evidentes de sobreajuste até que o
\textit{early stopping} interrompeu o treinamento. O DenseNet169 levou
mais tempo para estabilizar, mas manteve progressão consistente até
a época\,18. Já o ConvNeXt-Small atingiu seu melhor resultado na
época\,3 --- um comportamento notável. A combinação de
\textit{backbone} amplamente descongelado com imagens de 1024\,px parece
ter acelerado demais o aprendizado inicial, limitando a capacidade do
modelo de refinar suas representações nas épocas seguintes.

\begin{figure}[htbp]
\centering
\includegraphics[width=\columnwidth]{aprendizado.png}
\caption{Curva de aprendizado da loss do EfficientNetV2-S em 512\,px
(\textit{checkpoint} da época\,14) no conjunto de validação e treinamento.}
\label{fig:loss_v2}
\end{figure}

\begin{figure}[htbp]
\centering
\includegraphics[width=\columnwidth]{acuracia.png}
\caption{Curva de aprendizado da acurácia do EfficientNetV2-S em 512\,px
(\textit{checkpoint} da época\,14) no conjunto de validação e treinamento.}
\label{fig:acc_v2}
\end{figure}

\begin{figure}[htbp]
\centering
\includegraphics[width=\columnwidth]{aprendizado_densenet.png}
\caption{Curva de aprendizado da loss do DenseNet169
(\textit{checkpoint} da época\,18) no conjunto de validação e treinamento.}
\label{fig:loss_dense}
\end{figure}

\begin{figure}[htbp]
\centering
\includegraphics[width=\columnwidth]{acuracia_densenet.png}
\caption{Curva de aprendizado da acurácia do DenseNet169
(\textit{checkpoint} da época\,18) no conjunto de validação e treinamento.}
\label{fig:acc_dense}
\end{figure}

\begin{figure}[htbp]
\centering
\includegraphics[width=\columnwidth]{aprendizado_convnext.png}
\caption{Curva de aprendizado da loss do ConvNeXt-Small
(\textit{checkpoint} da época\,3) no conjunto de validação e treinamento.}
\label{fig:loss_dense}
\end{figure}

\begin{figure}[htbp]
\centering
\includegraphics[width=\columnwidth]{acuracia_convnext.png}
\caption{Curva de aprendizado da acurácia do ConvNeXt-Small
(\textit{checkpoint} da época\,3) no conjunto de validação e treinamento.}
\label{fig:acc_dense}
\end{figure}

\subsection{Matrizes de confusão}

\begin{figure}[htbp]
\centering
\includegraphics[width=\columnwidth]{matriz_confusao_effnet512.png}
\caption{Matriz de confusão do EfficientNetV2-S em 512\,px
(\textit{checkpoint} da época\,10) no conjunto de validação.}
\label{fig:confusao}
\end{figure}

A análise da matriz de confusão do EfficientNetV2-S em 512\,px
confirma que a classe \textit{Mild} é o ponto crítico do modelo.
Grande parte dos seus erros ocorre por confusão com \textit{No\,DR}
e \textit{Moderate}, o que faz sentido: lesões vasculares leves são
visualmente próximas tanto de retinas sem doença quanto de
estágios moderados da patologia. A classe \textit{Severe} também
apresentou confusões relevantes, principalmente com \textit{Moderate}
e \textit{Proliferative}. Em contrapartida, \textit{No\,DR} e
\textit{Moderate} tiveram as diagonais mais fortes. Vale notar que,
comparado aos experimentos em 384\,px, o modelo em 512\,px começou
a prever \textit{Severe} e \textit{Proliferative} de forma consistente,
o que não acontecia nas resoluções menores.

\subsection{Resultado do ensemble}

O ensemble por \textit{Weighted Soft Voting} apresentou desempenho
inferior ao esperado. Combinando seis modelos com pesos proporcionais ao
kappa individual de cada um, o resultado foi kappa de 0,4691 e F1 macro
de 0,39 --- pior que qualquer um dos três melhores modelos isolados,
como mostra a Tabela~\ref{tab:ensemble}.

\begin{table}[htbp]
\caption{Melhores modelos individuais versus ensemble}
\label{tab:ensemble}
\centering
\scriptsize
\begin{tabular}{l c c}
\hline
\textbf{Modelo} & \textbf{F1 Macro} & \textbf{Kappa} \\
\hline
ConvNeXt-Small (1024\,px)  & 0,53 & 0,75 \\
EfficientNetV2-S (512\,px) & 0,54 & 0,73 \\
DenseNet169 (512\,px)      & 0,55 & 0,71 \\
\textbf{Ensemble}          & \textbf{0,39} & \textbf{0,47} \\
\hline
\end{tabular}
\end{table}

A hipótese mais provável para esse resultado é a incompatibilidade
entre as distribuições de probabilidade dos modelos: quando as
saídas não estão calibradas numa mesma escala, a média ponderada
tende a produzir predições inconsistentes. Além disso, os pesos
do ensemble foram definidos de forma fixa, sem otimização sobre o
conjunto de validação. Não se pode descartar também algum
problema no carregamento dos \textit{checkpoints} de arquiteturas
distintas durante a execução. Assim, na configuração
avaliada, o ensemble não apresentou ganho suficiente para justificar
seu custo computacional adicional.

\subsection{Discussão dos principais achados}

Um dos achados mais relevantes ao longo dos experimentos foi o impacto
da resolução. Aumentar o EfficientNetV2-S de 384\,px para 512\,px
elevou o kappa de 0,57 para 0,73 e quase dobrou o recall de
\textit{Mild}, de 0,26 para 0,34. A explicação é direta:
microaneurismas e exsudatos característicos dos estágios iniciais da
retinopatia são estruturas pequenas, e resoluções menores podem
reduzir ou eliminar detalhes importantes que o modelo precisaria aprender
a reconhecer.

O desbalanceamento da base, por outro lado, provou ser mais resistente
do que o esperado. A ponderação de classes na função de
perda ajudou, mas não foi suficiente para compensar uma razão de
36:1 entre \textit{No\,DR} e \textit{Proliferative}. A situação
da classe \textit{Mild} é ainda mais delicada porque combina dois
problemas ao mesmo tempo: poucos exemplos e alta similaridade visual
com outras classes. O ConvNeXt-Small ilustra bem o risco de não
observar as métricas por classe: kappa de 0,7453, mas F1 de 0,09
para \textit{Mild}. Um modelo assim poderia passar por uma
avaliação superficial e ainda assim ser inapropriado para uso
clínico real.

Quanto ao ensemble, os resultados indicam que combinar modelos sem
calibração prévia pode ser contraproducente. Sem garantir que
as distribuições de probabilidade de cada modelo estejam numa
mesma escala, a combinação por média ponderada não agrega
informação de forma coerente. é uma lição importante:
ensemble não é uma solução automática para melhorar
desempenho --- a implementação importa tanto quanto a estratégia.

% =============================================================
\section{Conclusão}
% =============================================================

Este trabalho avaliou diferentes configurações de arquiteturas de
aprendizado profundo --- incluindo CNNs modernas e um Vision Transformer
--- para a classificação de cinco níveis de severidade de
retinopatia diabética em 35.126 imagens de fundo de retina. Foram
investigados o impacto da resolução de entrada, do fine-tuning
parcial do \textit{backbone}, da ponderação de classes e da
combinação de modelos por ensemble.

O EfficientNetV2-S em 512$\times$512 pixels apresentou o melhor
equilíbrio clínico, com kappa de 0,7313 e F1 macro de 0,54, sendo
um dos modelos com maior capacidade de detectar os cinco níveis de forma
consistente. O ConvNeXt-Small alcançou o maior kappa absoluto
(0,7453), mas seu recall de 0,05 para \textit{Mild} compromete sua
utilidade em cenários de triagem precoce. O ensemble não superou os
modelos individuais, evidenciando que a combinação de modelos
exige calibração adequada para ser efetiva.

De forma geral, os experimentos mostraram que resolução e
fine-tuning são fatores decisivos para o desempenho, mas que o
desbalanceamento extremo da base permanece o principal obstáculo.
Nenhuma das estratégias testadas foi suficiente para garantir boa
detecção em todas as classes simultaneamente --- especialmente em
\textit{Mild}, onde a dificuldade é tanto estatística quanto visual.

O estudo tem limitações relevantes: não foi utilizado conjunto
de teste externo independente, as restrições de memória de GPU
impediram explorar 1024\,px em todas as arquiteturas, e os experimentos
com \textit{focal loss} e \textit{sampler} não foram concluídos.
Para trabalhos futuros, as direções mais promissoras incluem
\textit{data augmentation} seletiva nas classes minoritárias ---
aplicando transformações mais intensas especificamente sobre
\textit{Mild}, \textit{Severe} e \textit{Proliferative} como forma de
oversampling sem introduzir ruído nas classes já bem representadas
---, validação cruzada estratificada, \textit{threshold moving}
por classe, calibração de probabilidades para viabilizar o ensemble
e aplicação de Grad-CAM para interpretar as regiões de interesse
identificadas pelos modelos.

% =============================================================
\begin{thebibliography}{00}
% =============================================================
\bibitem{grzybowski2022} A. Grzybowski, ``Artificial intelligence for
diabetic retinopathy screening,'' \textit{Acta Ophthalmol.}, vol.~100,
2022, doi: 10.1111/j.1755-3768.2022.15371.

\bibitem{gulshan2016} V. Gulshan et al., ``Development and Validation of
a Deep Learning Algorithm for Detection of Diabetic Retinopathy in
Retinal Fundus Photographs,'' \textit{JAMA}, vol.~316, no.~22,
pp.~2402--2410, Dec. 2016, doi: 10.1001/jama.2016.17216.

\bibitem{saygili2025} A. Sayg{\i}l{\i} and \"O.\,M. At{\i}c{\i},
``Enhancing Retinal Disease Detection With the Swin Transformer: A
Comprehensive Comparative Analysis,''
\textit{Concurrency Computat. Pract. Exper.}, vol.~37, no.~27--28,
p.~e70366, 2025, doi: 10.1002/cpe.70366.

\bibitem{rajeshwar2025} S. Rajeshwar, S. Thaplyal, A.\,M., and
G. Siva Shanmugam, ``Diabetic Retinopathy Detection Using DL-Based
Feature Extraction and a Hybrid Attention-Based Stacking Ensemble,''
\textit{Adv. Public Health}, vol.~2025, p.~8863096, 2025,
doi: 10.1155/adph/8863096.

\bibitem{dataset} ilovescience, ``Diabetic Retinopathy Resized,''
Kaggle, 2019. Disponível em:
https://www.kaggle.com/datasets/tanlikesmath/diabetic-retinopathy-resized.
Acesso em: 4 jun. 2026.

\bibitem{APTOS2019} Karthik, Maggie, e S. Dane, ``APTOS 2019 Blindness
Detection,'' Kaggle, 2019. Disponível em:
https://kaggle.com/competitions/aptos2019-blindness-detection.
Acesso em: 4 jun. 2026.
\end{thebibliography}

\end{document}